\section{Software Architecture}

The Siderust architecture separates reusable scientific mechanisms from
astronomy-specific interpretation. The lower layers represent quantities,
geometry and time; numerical components operate on these representations; and
the domain crate supplies astronomical models and composes complete workflows.
Figure~\ref{fig:architecture} illustrates this composition for a surface
observation. The separation gives each component a clear responsibility while
allowing scientific information to cross its interfaces.

\begin{figure*}[!tbp]
\centering
\begin{tikzpicture}[
  node distance=8mm,
  stage/.style={draw,rounded corners,align=center,text width=33mm,minimum height=13mm,font=\small},
  layer/.style={align=center,font=\footnotesize,text width=33mm},
  flow/.style={-{Latex[length=2mm]},thick}
]
\node[stage] (input) {Quantities and epoch\\Units, TT time, site};
\node[stage,right=of input] (reference) {Typed reference systems\\Center and frame};
\node[stage,right=of reference] (models) {Scientific models\\Ephemeris, orientation};
\node[stage,right=of models] (observable) {Observer geometry\\Altitude, azimuth, range};
\draw[flow] (input)--(reference);
\draw[flow] (reference)--(models);
\draw[flow] (models)--(observable);
\node[layer,below=3mm of input] {\texttt{qtty}, \texttt{tempoch}};
\node[layer,below=3mm of reference] {\texttt{affn}};
\node[layer,below=3mm of models] {\texttt{siderust-archive}\\Numerical components};
\node[layer,below=3mm of observable] {\texttt{siderust}\\Body-specific surface model};
\end{tikzpicture}

\caption{Scientific composition in a surface-observation workflow. Typed
quantities, epochs and reference systems connect ephemeris and orientation
models to observer-relative geometry. Reusable libraries supply the mechanisms;
the Siderust domain layer selects and combines the astronomical operations.
The diagram was prepared with AI-assisted TikZ authoring.}
\label{fig:architecture}
\end{figure*}

\subsection{Shared scientific representations}

\texttt{qtty} provides dimensional quantities with unit and scalar parameters.
Conversion changes the unit representation explicitly, while multiplication and
division construct compound quantities. \texttt{affn} provides generic Cartesian,
spherical and ellipsoidal containers, reference-frame and reference-center
traits, affine operators and conic geometry. \texttt{tempoch} supplies
scale-parameterized instants, durations, intervals and encoded epoch views.
Together, these libraries define the vocabulary shared by the scientific layers.

For example, a position uses quantities for its components and carries frame
and center markers through transformations. An epoch remains typed when passed
to an orientation model. The result of a local observation retains its observer
parameters while presenting altitude, azimuth and range as quantities.
This arrangement supports consistent reuse: the geometry library handles
coordinate structure, the time library handles temporal representation, and
the domain layer interprets both for a particular astronomical calculation.

\subsection{Models, numerical components and domain composition}

\texttt{siderust-archive} supplies scientific datasets, manifests, checksums
and generated accessors. Analytical ephemerides and time conversions consume
these data through their respective implementations. \texttt{cheby} supplies
Chebyshev approximation and interpolation; \texttt{keplerian} supplies
central-force mechanics, conic propagation and Lambert solvers;
\texttt{principia} supplies numerical dynamics, integration and variational
equations; and \texttt{optica} supplies optical and radiometric primitives.
These components can be reused independently or assembled by the domain crate
(Table~\ref{tab:crate-products}).

The \texttt{siderust} crate adds celestial bodies, observatories, astronomical
transformations, analytical and runtime ephemerides, observing-event searches
and feature-gated orbital workflows. Generic surface observers illustrate
this division particularly well: a common projection pipeline operates on a
body-specific reference surface and orientation provider. Extending the
observer model therefore reuses the quantity, epoch and geometry layers
alongside the local-horizon calculation.
This separates two extension tasks: a new surface model supplies scientific
parameters, while the existing projection consumes the shared typed geometry.
The boundary cases described in Section~\ref{sec:evaluation} examine this
separation without changing the underlying orientation or ephemeris implementation.

\begin{table}[!tbp]
\centering
\caption{Reusable components and their role in scientific composition. Exact
versions and licensing terms are available from the published software releases.}
\label{tab:crate-products}
\footnotesize
\begin{tabularx}{\columnwidth}{@{}lY@{}}
\toprule
Component & Scientific responsibility \\
\midrule
\texttt{qtty} & Quantities, units and dimensional arithmetic \\
\texttt{affn} & Geometry, coordinate containers, frames and centers \\
\texttt{tempoch} & Epochs, time scales, durations and conversions \\
\texttt{siderust-archive} & Scientific datasets and provenance \\
\texttt{cheby} & Approximation, interpolation and calculus \\
\texttt{keplerian} & Analytic orbital mechanics and transfers \\
\texttt{principia} & Numerical dynamics and state transport \\
\texttt{optica} & Optical and radiometric operations \\
\texttt{siderust} & Astronomical models and integrated workflows \\
\bottomrule
\end{tabularx}
\end{table}

\subsection{Memory safety and scientific semantics}

Rust's ownership, borrowing and lifetime rules provide the memory and sharing
foundation for these components. Safe Rust prevents important classes of
invalid memory access and data races, assuming sound implementations of the
unsafe abstractions on which it relies \cite{rustbook,rustbelt2018}.
Siderust adds domain-level structure: unit, scale, frame and center types
constrain scientific operations, while fallible constructors and explicit
errors handle value-dependent conditions.

The released domain crate inherits a workspace
\texttt{unsafe\_code = "forbid"} lint. Its FFI package has a separate lint
configuration allowing unsafe code at the foreign boundary; dependencies have
their own implementation policies. This source policy complements API design
and testing. Scientific validity additionally depends on choosing suitable
models, data and inputs, as discussed in Section~9.
